\begin{itemize}
\item \textbf{Retrieval tokens.} The handoff packet is counted three ways:
\begin{itemize}
\item exact UTF-8 bytes;
\item a named tokenizer in MCP-03 (Qwen2.5-0.5B-Instruct, pinned revision);
\item a labelled 4-bytes-per-token estimate in the MCP-02 economics.
\end{itemize}
The tokens the router \emph{read} locally (every tool result) are counted separately, as are router tokens (encoder
positions processed).
\item \textbf{Compute.} FLOPs per candidate come from \texttt{torch.utils.flop\_counter.FlopCounterMode} at the fixed sequence
length, multiplied by the candidates scored. Recurrent depth is the number of block applications. Latency is the
wall clock per router decision on the recorded host.
\item \textbf{Memory.} For training we record the process peak working set, which includes the in-memory dataset, and the CUDA
peak allocation on GPU. For inference we record the peak working-set delta around one 32-candidate forward pass in a
fresh process, together with the weight bytes.
\item \textbf{Economics.} Cost per accepted task divides the summed cost of every assigned episode, failures included, by
the accepted count. The downstream handoff is priced at an assumed frontier input price, and local seconds at
amortised device prices. Tool calls and verifier checks carry nominal prices. All prices are labelled assumptions,
with four sensitivity scenarios.
\end{itemize}
The hosts were shared: an RTX 5080 Laptop GPU with other workloads resident (80 W limit), and a CPU at
full load from concurrent jobs. Wall-clock figures are upper bounds; FLOPs and token counts are host independent.
